\section{Schlussfolgerung und Ausblick}

\subsection{Zusammenfassung}

MapApp zeigt, dass die in M1 identifizierte Lücke im Bereich mobiler
Open-Data-Anwendungen für Wien gezielt geschlossen werden kann.
Von $n = 8$ bestehenden mobilen Apps auf \textit{data.gv.at} war keine
auf die allgemeine, kartenbasierte Exploration von Open Data ausgerichtet.
MapApp adressiert genau dieses Desiderat und verbindet in einer kohärenten
Material-Design-3-Oberfläche drei Funktionsbereiche:
interaktive Kartenexploration, Datensatzverwaltung und Personalisierung.

\textbf{Stärken:}
\begin{itemize}
  \item Die farbkodierte Pindarstellung – motiviert durch \cite{Visualization} –
        ermöglicht sofortige kategoriale Orientierung ohne Legende.
  \item Die bidirektionale Navigation zwischen Karte und Downloads
        erfüllt alle vier Nutzeraufgaben aus M1 (\textit{Erkunden},
        \textit{Suchen}, \textit{Teilen}, \textit{Schnell informieren})
        ohne Bildschirmwechsel-Hürden.
  \item Das Barrierefreiheitssystem (Großschrift, Hochkontrast, Systemdesign)
        richtet sich an die heterogenen Nutzergruppen, die Simonofski et al.\
        \cite{simonofski2022} als zentrale Designanforderung identifizieren.
  \item Offline-Download löst das in M1 als kritisch eingestufte Problem
        langer Ladezeiten bei schwacher Verbindung strukturell und proaktiv.
\end{itemize}

\textbf{Schwächen und Einschränkungen:}
\begin{itemize}
  \item Die Datenbasis ist statisch im Arbeitsspeicher gehalten.
        Eine Live-Anbindung an \textit{data.wien.gv.at} \cite{datawien}
        überstieg den Projektrahmen, ist aber für eine produktive Version
        unerlässlich.
  \item Der Download-Fortschritt ist simuliert; echte Offline-Kacheln
        werden nicht persistent gespeichert.
  \item Die Pin-Klickerkennung basiert auf euklidischer Näherung statt
        echtem Layer-Hit-Testing, was bei eng beieinanderliegenden Pins
        zu Fehlauswahlen führen kann.
\end{itemize}

\subsection{Interessante nächste Schritte}

\begin{itemize}
  \item \textbf{Erweiterter Datensatzkatalog:}
        Einbindung weiterer Wiener Open-Data-Themen wie Luftqualität,
        Mietpreise oder Bevölkerungsdaten – damit auch die Personas Ralph
        (beruflich) und Benjamin (Forschung) vollständig bedient werden.
  \item \textbf{Filter- und Kategorisierungsfunktionen:}
        In M2 wurden auf Basis von Nutzerfeedback erweiterte Suchfilter
        (reguläre Ausdrücke, Lizenzfilter) als wünschenswert identifiziert.
        Diese könnten in einer Folgeversion umgesetzt werden.
  \item \textbf{Routenplanung:}
        Integration einer Routenführung, die Open-Data-Ebenen
        (z.\,B.\ Radwege, ÖPNV-Stationen) direkt in die Berechnung einbezieht.
\end{itemize}
